\documentclass[a4paper, 12pt]{article}
\usepackage{amsmath, amssymb, amsthm}
\usepackage[utf8]{inputenc}
\usepackage[ngerman]{babel}
\usepackage[backend=biber]{biblatex}
\usepackage{enumitem}
\usepackage[left=3cm, right=2cm]{geometry}
\usepackage[intlimits]{esint}
%\usepackage{tikzit}
\usepackage{tikz}
\setlength{\parindent}{0pt}\begin{document}
%\input{sample.tikzstyles}
\usetikzlibrary{decorations, decorations.pathreplacing, decorations.pathmorphing}
\newcommand{\re}{\text{Re\quad}}
\newcommand{\im}{\text{Im\quad}}
%A tikz picture as an equation:
%\ctikzfig{tikzfig}

\subsection{Budyko model}
\begin{tikzpicture}
\def\r{9}
\def\t{20}
\def\p{35}
\def\w{0.5mm}
\draw[line width = \w] (\r,0) arc[start angle = 0, end angle=90, radius = \r];

\draw[->, dashed, line width = \w] (0,\r) -- (0,\r + 1) node[above]{\LARGE N};

\draw[dashed, line width = \w] (0,0) -- ({\r * cos(5)},{\r * sin( 5 ) });
\draw[dashed, line width = \w] (0,0) -- ({\r * cos(\t)},{\r * sin( \t ) });
\draw[dashed, line width = \w] (0,0) -- ({\r * cos(\p)},{\r * sin( \p ) });

%\draw[line width = 0.75*\w, ->] (2,0)  arc[start angle = 0, end angle = \t, radius = 2];
%\draw[line width = 0.75*\w, ->] (3,0)   arc[start angle = 0, end angle = \p, radius = 3];

%\draw ({1 * cos(\t/2)},{1* sin( \t/2 ) }) node[left]{\Large $\theta_1$};

\draw[->, line width=\w, red ] ({\r * cos(\t /2 + 3) + 4},{\r * sin( \t /2 +3) }) node[right=1mm, black]{\LARGE $S_{sol}(\theta_1)$} -- ({\r * cos(\t /2 + 3) },{\r * sin( \t /2 +3 ) });
\draw[->,line width=0.5*\w, red] ({\r * cos(\t /2 + 3) },{\r * sin( \t /2 + 3) }) -- ({\r * cos(\t /2+ 3) + 4},{\r * sin( \t /2 + 3) +1}) node[right=1mm, black] {\LARGE $\alpha(\theta_1)S_{sol}(\theta_1)$};

\draw[->, decoration={snake}, decorate, red, line width=\w] ({\r * cos(\t /2 - 2) },{\r * sin( \t /2 - 2) }) -- ({\r * cos(\t /2 -2 ) + 4.05},{\r * sin( \t /2 -2)}) node[left=10mm, black, below=2mm] {\Large $S_{OLW}\sim A + B(T(\theta_1))$};

\draw[red, ->, line width = \w] ({\r * cos(\t -5)},{\r * sin( \t -5 ) })  arc[start angle = \t-5, end angle=\t+5 , radius = \r];

\draw ({\r * cos(\t)},{\r * sin( \t ) }) node[right,rotate= \t]{\Large Heat transport};

%\draw ({\r * cos(\t/2-4)},{\r * sin( \t/2-4) }) node[left=25mm, above= -3mm, rotate=0]{\Large with $\alpha(\theta_1), \, T(\theta_1)$};

\draw ({\r * cos(\t/2+1.5)},{\r * sin( \t/2) }) node[left, rotate=0]{\Large Latitude $\theta_1$};

\draw ({\r * cos(\p/2+5)},{\r * sin( \p/2+5) }) node[left, rotate=0]{\Large Latitude $\theta_2$};

\end{tikzpicture}


\subsection{Conservation Bathtub}


\begin{tikzpicture}
			\draw [thick] (3,3) ellipse (3 and 0.5);
			\draw [thick] (0,3) to [out=270,in=180] (2,0);
			\draw [thick] (2,0) -- (4,0);
			\draw [thick] (4,0) to [out=0,in=270] (6,3);
			\draw [thick] (6,3) -- (6,5) -- (4,5) -- (4,4.5);
			\draw [thick,cyan] (3,1.6) ellipse (2.74 and 0.4);
			\draw [->,thick, green,dashed] (4,4.5) -- (4,1);
			\draw [->,thick,red,dashed] (4.5,0.5) -- (8,0.5);
			\draw [thick, green] (7,5) -- (7.5,5);
			\node [right] at (7.5,5) { Water in};
			\draw [thick, red] (7,4) -- (7.5,4);
			\node [right] at (7.5,4) { Water out};
\end{tikzpicture}

\subsection{Domains}
\begin{tikzpicture}
\draw (0,0) ellipse (2cm and 1cm);
\draw[ <-] (0,0) -- (2,1) node[right]{$\omega$};
\draw[ <-] (2,0) -- (3,1) node[right]{$\partial\omega$};
\draw[ ->] (-2,0) -- (-3,0) node[above]{$\vec{n}(x)$};
\draw[fill=black] (-2,0) circle (0.05) node[right]{$x$};

\end{tikzpicture}

\subsection{Time discretization}
%\usepgfmodule{decorations}
\begin{tikzpicture}
\draw[thick,->] (-1,0) -- (9,0) node[right] {$t$};
\draw[thick,->] (-1,0) -- (-1,3) node[above] {$y(t)$};
\foreach \x in {0,1,2,3}
   \draw (2*\x ,2pt) -- (2*\x,-2pt) node[below] {$t_{\x}$};
\draw (8 ,2pt) -- (8 ,-2pt) node[below] {$\cdots$};

\draw [decorate,decoration={brace,amplitude=10pt},xshift=0pt,yshift=0pt]
(0 ,0.6 ) -- (2,0.6 ) node [black,midway,yshift=17pt] 
{$\Delta t$};
\end{tikzpicture}


\subsection{Integral approximation}
\begin{tikzpicture}
[domain=-1:6.1]
\draw[thick,->] (-1,0) -- (9,0) node[right] {$t$};
\draw[thick,->] (-1,0) -- (-1,7) node[above] {$f(y,t)$};
\foreach \x in {0,1,2,3}
   \draw (2*\x ,2pt) -- (2*\x,-2pt) node[below] {$t_{\x}$};
\draw (8 ,2pt) -- (8 ,-2pt) node[below] {$\cdots$};
\draw[color=red] plot(\x,{(-3/16)*\x*\x*\x+(11/8)*\x * \x -(3/2)*\x +2});
\foreach \x in {0,2,4}
   \draw [color=blue](\x ,{(-3/16)*\x*\x*\x+(11/8)*\x * \x -(3/2)*\x +2}) -- (\x + 2,{(-3/16)*\x*\x*\x+(11/8)*\x * \x -(3/2)*\x +2});
\foreach \x in {0,2,4}
   \draw[color=blue, dotted] (\x ,{(-3/16)*\x*\x*\x+(11/8)*\x * \x -(3/2)*\x +2}) -- (\x , 0);
\draw[color=blue, dotted] (6,6)--(6,0);
\foreach \x in {0,2,4}
   \draw [fill=blue,blue, draw opacity=0.2, fill opacity=0.2](\x ,{(-3/16)*\x*\x*\x+(11/8)*\x * \x -(3/2)*\x +2}) rectangle (\x + 2,0);

\end{tikzpicture}
\subsection{Realteil und Imaginärteil}
\begin{tikzpicture}
[domain=0:6]
\draw[thick,->] (0,0) -- (9,0) node[right] {$t$};
\draw[thick,->] (0,-3) -- (0,3) node[above] {Im$(y(t))$};


\draw[color=red] plot[samples=150](\x,{exp(0.2*\x)*sin(10*\x r)}) node[right] {$e^{at}\sin(bt)=\text{Im}(y(t))$};
\draw[color=blue, dotted] plot[samples=150](\x,{exp(0.2*\x)}) node[right] {$e^{at}$};
\draw[color=blue, dotted] plot[samples=150](\x,{-exp(0.2*\x)}) node[right] {$-e^{at}$};
\node[draw,color=violet] at (2,2.5) {$a>0$};

\end{tikzpicture}
\\
\begin{tikzpicture}
[domain=0:6]
\draw[thick,->] (0,0) -- (9,0) node[right] {$t$};
\draw[thick,->] (0,-3) -- (0,3) node[above] {Im$(y(t))$};


\draw[color=red] plot[samples=150](\x,{exp(-0.2*\x)*sin(10*\x r)}) node[above, yshift=20 pt] {$e^{at}\sin(bt)=\text{Im}(y(t))$};
\draw[color=blue, dotted] plot[samples=150](\x,{exp(-0.2*\x)}) node[right] {$e^{at}$};
\draw[color=blue, dotted] plot[samples=150](\x,{-exp(-0.2*\x)}) node[right] {$-e^{at}$};

\node[draw,color=violet] at (2,2.5) {$a<0$};


\end{tikzpicture}

\subsection{Circles}
\begin{tikzpicture}
\draw[thick,->] (-3,0) -- (3,0) node[right] {$\text{Re}(\lambda)\cdot\Delta t$};
\draw[thick,->] (0,-3) -- (0,3) node[above] {$\text{Im}(\lambda)\cdot\Delta t$};

\draw[fill = red,red,fill opacity=0.2] (-1,0) circle (1);

\draw (-2 ,2pt) -- (-2,-2pt) node[anchor=north east] {$-2$};
\node at (0.2,-0.3) {$0$};
\draw[color=red, dotted, line width=0.7pt](-1,1)--(0,1);
\draw[color=red, dotted, line width=0.7pt](-1,-1)--(0,-1);
\draw (-2pt ,1) -- (2pt,1) node[right] {$1$};
\draw (-2pt ,-1) -- (2pt,-1) node[right] {$-1$};
\node[draw,color=violet] at (-1,-2) {$a<0$};
\node[draw,color=violet] at (1,-2) {$a>0$};


\end{tikzpicture}
\\
\begin{tikzpicture}
\draw[thick,->] (-3,0) -- (3,0) node[right] {$\text{Re}(\lambda)\cdot\Delta t$};
\draw[thick,->] (0,-3) -- (0,3) node[above] {$\text{Im}(\lambda)\cdot\Delta t$};

\draw[red, dotted, line width=0.7pt] (1,0) circle (1);
\def\firstcircle{(1,0) circle (1)};
\draw (2 ,2pt) -- (2,-2pt) node[anchor=north west] {$2$};
\node at (-0.2,-0.3) {$0$};

\draw (-2pt ,1) -- (2pt,1) node[left] {$1$};
\draw (-2pt ,-1) -- (2pt,-1) node[left] {$-1$};
%\draw[fill=red,red, draw opacity=0.2, fill opacity=0.2] (3,3) rectangle (-3,-3);
\begin{scope}[even odd rule]% first circle without the second
            \clip \firstcircle (-3,-3) rectangle (3,3);
        \fill[red, fill opacity=0.2] (3,3) rectangle (-3,-3);
\end{scope}

\end{tikzpicture}

\subsection{Spatial discretization}
%\usepgfmodule{decorations}
\begin{tikzpicture}
\draw[thick,->] (-1,0) -- (9,0) node[right] {$x$};
\draw[thick,->] (-1,0) -- (-1,4) node[above] {$T(x,t)$};
\draw (2*0 ,2pt) -- (2*0,-2pt) node[below] {$\cdots$};
\draw (2*1 ,2pt) -- (2*1,-2pt) node[below] {$x_{i-1}$};
\draw (2*2 ,2pt) -- (2*2,-2pt) node[below] {$x_{i}$};
\draw (2*3 ,2pt) -- (2*3,-2pt) node[below] {$x_{i+1}$};
\draw (8 ,2pt) -- (8 ,-2pt) node[below] {$\cdots$};
\draw[color=red,domain=-0.5:8.5] plot[samples=92,mark=x,mark indices={27,47,67}](\x,{0.1*((-2/16)*\x*\x*\x+(11/8)*\x * \x -(3/2)*\x)+2});

\node at (2,2.5) {$T_{i-1}$};
\node at (4,3.2) {$T_{i}$};
\node at (6,3.7) {$T_{i+1}$};

\draw [decorate,decoration={brace,amplitude=10pt},xshift=0pt,yshift=0pt]
(0 ,0.6 ) -- (2,0.6 ) node [black,midway,yshift=17pt] 
{$\Delta x$};
\end{tikzpicture}

\end{document}
